\section{Z. Additional steps in cote 149, scans 1--40}

\begin{result}{B11}{Refuted}{An affine hyperbolic polycurve need not immerse in a semiabelian variety}
\label{res:B11}
The semiabelian-immersion assertion on scan 26, repeated parenthetically
on scan 32, is false for the class of successive elementary hyperbolic
fibrations used in Corollary 1(b). There is a smooth affine hyperbolic
polycurve of dimension two every morphism from which to a semiabelian
variety contracts its curve fibres; this occurs even when both levels
have genus zero. This gives sweep rows ZAB1 and ZAB6 (149, section 1,
manuscript pp.~12, 15 and 17, scans 26, 32 and 36), including the
parenthetical claim that genus-zero fibres supply a global torus immersion.
B5 proves the stated corollaries without this assertion.
\end{result}
\begin{proof}
Work over $k=\overline{\mathbb Q}$. Let
\[
 B:\quad v^2=u(u-1)(u-2)(u-3)(u-4)
\]
be its smooth projective genus-two model. Let $P_0,P_1$ be the
Weierstrass points over $0,1$. The function $u/(u-1)$ has divisor
$2P_0-2P_1$ and is not a square in $k(B)$: a square root would have
divisor $P_0-P_1$, giving a degree-one function on a genus-two curve.
Let $C$ be the normalization of $B$ in
$k(B)(\sqrt{u/(u-1)})$. All valuations of this function are even,
so this connected double cover is \'etale: locally remove the even
uniformizer power and take a square root of the resulting unit.
Write $\iota_C$ for its fixed-point-free deck involution.

Let $D$ be another copy of $B$, let $\iota_D$ be its hyperelliptic
involution, and choose a fixed Weierstrass point $W$. Remove any point
$b_0$ from $B$, its two preimages from $C$, and $W$ from $D$;
write the resulting affine curves as $B^\circ,C^\circ,D^\circ$.
The diagonal involution on $C^\circ\times D^\circ$ is free, and
\[
 Y=(C^\circ\times D^\circ)/\langle\iota_C\times\iota_D\rangle
         \longrightarrow B^\circ
\]
is a smooth affine fibration with fibre $D^\circ$. The quotient is
affine because it is a finite-group quotient of an affine scheme;
it is smooth because the free quotient map is \'etale. The family
becomes the product after the \'etale cover $C^\circ\to B^\circ$.
Its smooth proper compactification with boundary descends from
$C^\circ\times D$ and the invariant section $W$. Both base and
fibre have hyperbolic type $(2,1)$.

Given $f:Y\to G$ with $1\to T\to G\to A\to1$ semiabelian,
pull it back to $h:C^\circ\times D^\circ\to G$.
The composite into $A$ extends to $C\times D$ by the opened
\cite[I, Theorem 3.2]{AB_MilneAV}: rational maps from nonsingular
varieties to abelian varieties extend everywhere. By
\cite[I, Corollary 1.5]{AB_MilneAV}, this extension has the form
$a(c)+b(d)+a_0$. Diagonal invariance implies that
$b(\iota_Dd)-b(d)$ is constant in $d$. After choosing $W$ as origin
for the Abel map, $b-b(W)$ is induced by a homomorphism
$\beta:J_D\to A$. The hyperelliptic involution acts as $-1$ on
$J_D$, because $d+\iota_Dd$ is linearly equivalent to $2W$.
Thus $-2\beta$ is a constant homomorphism, so $\beta=0$ (its
connected image would otherwise be killed by two). The map into
$A$ is independent of $d$.

For every fixed $c$, the remaining map $D^\circ\to G$ lies in a
torus fibre and is constant. Indeed, every unit on $D\setminus\{W\}$
is constant: its divisor is supported at $W$ and has degree zero.
The same is true coordinate by coordinate after identifying the
fibre with $\mathbb G_m^{\dim T}$. Hence $f$ contracts every curve
fibre of $Y\to B^\circ$, and cannot be an immersion. The script
checks the five branch values are distinct; the divisor and
ramification verifications above are the complete algebraic checks.

Here is the genus-zero example needed to check the additional assertion
on scan 36. Let $\zeta$ be a primitive cube root of unity and put
\[
 B_0=\mathbb P^1\setminus\{0,1,\infty\},\quad
 C_0=\mathbb P^1\setminus\{0,\infty,1,\zeta,\zeta^2\},\quad
 D_0=\mathbb P^1\setminus\{1,\zeta,\zeta^2\}.
\]
The map $C_0\to B_0$, $z\mapsto z^3$, is finite \'etale Galois
of degree three. On $C_0\times D_0$ let the generator act by
$(z,w)\mapsto(\zeta z,\zeta w)$, and set
\[
             Y_0=(C_0\times D_0)/\langle\zeta\rangle.
\]
As above this is a smooth affine polycurve over $B_0$, now with
both base and geometric fibre of type $(0,3)$. The boundary of its
relative proper compactification is finite \'etale of degree three;
it need not consist of three globally labelled sections.

Every invertible regular function on $C_0\times D_0$ has the form
$f(z)g(w)$. To prove this, consider its divisor on $C_0\times\mathbb P^1$.
It is supported on the three irreducible divisors $C_0\times\{\zeta^i\}$.
A restriction at a suitable $z_0$ has the same three orders. Dividing
by that restriction removes the divisor, and the resulting invertible
function on $C_0\times\mathbb P^1$ comes from $C_0$, by properness
of the second factor. An invariant function $f(z)g(w)$ satisfies
$g(\zeta w)/g(w)=f(z)/f(\zeta z)$, a constant. Thus the divisor
of $g$ is invariant under the cyclic permutation of the three cusps.
Its three coefficients are equal, and their sum is zero; all vanish.
Therefore $g$ is constant. All units on $Y_0$ are constant on the fibres
of $Y_0\to B_0$.

For a map $Y_0\to G\to A$, its pullback extends to
$\mathbb P^1\times\mathbb P^1\to A$ by the same opened extension
theorem. It is constant on each projective-line factor, since a
projective line has zero Jacobian, and hence is constant. Translating
the image fibre of $G\to A$ identifies the remaining map with a map
to a torus. Its coordinates are the units just computed, so it too
contracts every curve fibre. This proves the asserted failure of a
global torus or semiabelian immersion even for genus-zero towers.
\end{proof}

\begin{result}{A7}{Proved}{The field-map image, conservativity, and the limit groupoid}
\label{res:A7}
For a finitely generated field extension $L/F$, restriction
$G_L\to G_F$ has open image $G_{F_s}$, where $F_s$ is the relative
separable algebraic closure of $F$ in $L$. The map is injective
exactly when $L/F$ is algebraic, and is an isomorphism exactly when
$L/F$ is finite purely inseparable. Thus the field functor of
scans 5--6 is conservative in characteristic zero, and becomes
conservative after inverting finite purely inseparable extensions
in positive characteristic.
For an arbitrary field $F=\varinjlim F_i$, where $F_i$ runs through
its finitely generated prime-field subfields, the compatible
absolute groups satisfy $G_F\simeq\varprojlim G_{F_i}$; the analogous
statements hold for the profinite groupoids and their finite-cover
topoi. Keeping the whole pro-system retains the faithfulness of
B10, whereas taking only its limit does not. This closes sweep
row ZAB2 (149, section 1, manuscript pp.~2--3, scans 5--8).
\end{result}
\begin{proof}
Choose one algebraic closure of $L$ and the compatible separable
closures. The fixed field in $F^{\rm sep}$ of the restriction image
is $L\cap F^{\rm sep}=F_s$; compactness makes the image closed,
and Galois correspondence identifies it with $G_{F_s}$. The field
$F_s/F$ is finite because $L/F$ is finitely generated. The kernel
is the absolute group of $LF^{\rm sep}$, up to its purely
inseparable closure, which leaves finite \'etale extensions unchanged.
If $L/F$ is algebraic, this kernel is one. If its transcendence
degree is positive, choose a divisorial discrete valuation of the
finitely generated field $LF^{\rm sep}/F^{\rm sep}$. A uniformizer
has a nontrivial separable $\ell$th-root extension for
$\ell\ne\operatorname{char}F$, since its value is one. The kernel
is nontrivial. Injectivity plus surjectivity is thus equivalent to
algebraicity and $F_s=F$, or to finite pure inseparability.

Inside $F^{\rm sep}$ choose the separable algebraic closure of each
$F_i$. Their directed union is $F^{\rm sep}$: a separable polynomial
and a witness to its nonzero discriminant use finitely many
coefficients, as does every finite collection of such polynomials.
Compatible automorphisms of these subfields therefore extend uniquely
to $F^{\rm sep}$ and fix $F$, giving the group isomorphism. Conversely
restriction supplies every component. It is a homeomorphism by
compactness. A finite \'etale algebra and all its morphisms descend
to a sufficiently large $F_i$ by their finite presentations; this
proves the groupoid and topos comparisons as well.

Finally, $\operatorname{Ind}(\mathrm{Fields}_{\rm fg})\to
\mathrm{Fields}$ is an equivalence: every field is such a union,
and a map from a finitely generated field into a directed union
lands in one stage, with equality detectable at a common later
stage. For a faithful contravariant functor on the finitely generated
objects, the formula for pro-morphisms,
$\operatorname{Hom}_{\rm Pro}(X_i,Y_j)=
\varprojlim_j\varinjlim_i\operatorname{Hom}(X_i,Y_j)$,
therefore preserves injectivity on these field maps; equality at a
stage in the target can be pulled back by faithfulness. We use the
relative identity marking of B10. The limit alone can lose this:
separably closed fields have trivial absolute group, while their
field embeddings can be distinct. For instance $\overline{\mathbb Q}$
has distinct automorphisms but its absolute group is one.
\end{proof}

\begin{result}{A8}{Proved}{Cofinal Artin models and the dimension of their groups}
\label{res:A8}
In characteristic zero every smooth integral variety has a dense
open which is a tower of affine smooth hyperbolic curves with
smooth proper compactifications and finite \'etale boundary.
Thus such opens are cofinal among models of a finitely generated
field, and their full geometric groups are successive extensions
of nonabelian free profinite groups. A tower of length $n$ has
$\ell$-cohomological dimension exactly $n$ for every prime.
This closes the additional assertions on scans 10--13, sweep
row ZAB3 (149, section 1, manuscript pp.~4--6).
\end{result}
\begin{proof}
Let $V/k$ be smooth integral of dimension $n$. Choose a separating
transcendence basis $t_1,\ldots,t_n$ of $k(V)/k$, and project on
$t_1,\ldots,t_{n-1}$. Over $F=k(t_1,\ldots,t_{n-1})$ the generic
fibre has a smooth projective curve model: normalize $\mathbb P^1_F$
in the finite extension $k(V)/F(t_n)$. Characteristic zero makes
this regular curve smooth. This model and its birational comparison
to the generic fibre spread over a nonempty open
$B\subset\mathbb A^{n-1}_k$. Shrinking $B$ makes the model smooth
proper with geometrically connected fibres; regularity of the
function-field extension gives geometric connectedness. If the
chosen projection initially has a larger constant field in its
generic fibre, replace $B$ by its normalization in that finite
relative algebraic closure; after shrinking it is smooth and the
same construction applies over its function field.

The complement of the comparison open in the generic proper curve
is finite. Add finitely many distinct separable closed points so
that the complement is nonempty and hyperbolic. Spread these
points to disjoint finite \'etale boundary divisors after shrinking
$B$. Remove any remaining vertical bad locus. Their complement
$U$ is a smooth curve fibration contained in $V$. It is affine over
$B$, since a positive-degree relative divisor on a proper smooth
curve is relatively ample. Apply the same construction inductively
to a dense affine open of the smooth base $B$ and pull $U$ back to
it. Each comparison is the actual restriction to a dense open,
so the final tower lies inside the originally prescribed $V$.
Applying this to every dense open proves cofinality. C4 gives the
exact sequences; complex comparison and A3 identify the affine
hyperbolic fibre groups as finite-rank free profinite groups.

Here is a proof of the exact dimension, not just its upper bound.
A free profinite group of positive finite rank has cohomological
dimension one: generators lift across every finite surjection,
so every finite embedding problem is solvable; equivalently all
$H^2$ with finite modules vanish, and dimension shifting gives the
upper bound. Its nonzero $H^1(-,\mathbb F_\ell)$ gives equality.
Hochschild--Serre gives the additive upper bound for a tower.
For the reverse inequality in $1\to N\to P\to Q\to1$, where
$N$ is free of rank $r>0$ and $\operatorname{cd}_\ell Q=n-1$,
first take an open subgroup of $Q$ acting trivially on
$H^1(N,\mathbb F_\ell)$. Its dimension is still $n-1$ by Serre's
finite-index theorem: successive extensions of the torsion-free
groups in A3 are torsion-free.

Some further open subgroup $Q'$ has
$H^{n-1}(Q',\mathbb F_\ell)\ne0$. Indeed choose a finite
$\mathbb F_\ell[Q]$-module $M$ with nonzero top cohomology and a
surjection onto $M$ from finitely many copies of a permutation
module $\mathbb F_\ell[Q/Q']$. Vanishing in degree $n$ for its
kernel makes the top-cohomology map surjective, and Shapiro's
lemma gives the assertion. For the preimage $P'$ of $Q'$, the
Hochschild--Serre term
\[
 E_2^{n-1,1}=H^{n-1}(Q',\mathbb F_\ell)^{r}
\]
is nonzero. No differential enters it (the fibre has no degree-two
cohomology) or leaves it (the base has dimension $n-1$).
Thus $H^n(P',\mathbb F_\ell)\ne0$, proving the lower bound.
For $n=1$ use the initial free group case. This proves the
manuscript's finite-model equality as well as the limiting field
calculation A1.
\end{proof}

\begin{result}{A9}{Proved}{Reconstruction of a centerless extension and uniqueness of lifted maps}
\label{res:A9}
For a profinite extension $1\to P\to E\to\Gamma\to1$ with
$Z(P)=1$, the outer action determines the extension, with its
identified kernel and quotient, by
\[
 E\ \xrightarrow{\sim}\
 \operatorname{Aut}(P)\times_{\operatorname{Out}(P)}\Gamma,
 \qquad e\longmapsto(\operatorname{int}(e)|_P,\bar e).
\]
Use the compact-open topology on continuous automorphisms and the
quotient topology on outer automorphisms; the displayed pullback
is profinite because it is isomorphic to $E$.
If two extension maps over $\Gamma$ restrict to the same map
$P\to P'$, and its image has trivial centralizer in $P'$, the
extension maps are equal. Conversely a continuous homomorphism
$\varphi:P\to P'$ whose class modulo target inner automorphisms
commutes with the two outer $\Gamma$-actions, and whose image has
trivial centralizer in $P'$, extends uniquely to a continuous
homomorphism $E\to E'$ over $\Gamma$. These are the formal claims of scans
16--20 and 29, sweep row ZAB4 (149, section 1, manuscript pp.~7--9,
14). They do not require $\operatorname{Aut}(P)$ itself to be
profinite for an arbitrary infinitely generated $P$.
\end{result}
\begin{proof}
The displayed homomorphism is injective because its kernel is
$Z(P)$. For a pair $(a,\gamma)$ in the pullback, choose $e$ above
$\gamma$. The equality of outer classes means
$a\operatorname{int}(e)^{-1}=\operatorname{int}(p)$ for a unique
$p\in P$. Thus $pe$ maps to the pair. Conjugation by elements of
$E$ varies continuously in the compact-open topology, and
$\operatorname{Inn}(P)$ is the compact, hence closed, image of $P$.
The continuous bijection from compact $E$ to this Hausdorff
pullback is a homeomorphism. This verifies both the algebra and
topology of the reconstruction.
For extension maps $F,G$, the element $G(e)F(e)^{-1}\in P'$
centralizes their common image of $P$, by evaluating the two maps
on $epe^{-1}$. The centralizer hypothesis makes it one for every
$e$, proving uniqueness. An outer equality is handled first by
the specified inner conjugation, exactly as in B5.
For existence, choose $e\in E$ with image $\gamma\in\Gamma$, and
any lift $e'_0\in E'$ of $\gamma$. Outer equivariance says exactly
that some $q\in P'$ makes $e'=qe'_0$ satisfy
\[
 \operatorname{int}(e')\circ\varphi
       =\varphi\circ\operatorname{int}(e)|_P.
\]
The centralizer hypothesis makes $e'$ unique. Define $F(e)=e'$.
Uniqueness proves $F(e_1e_2)=F(e_1)F(e_2)$ and
$F|_P=\varphi$. The graph of $F$ is the closed subset of
$E\times E'$ defined by equality of quotient images and the
displayed equations for all $p\in P$. Its projection to $E$ is a
continuous bijection from a compact space to a Hausdorff space,
hence a homeomorphism; $F$ is continuous.

Finally choose an object in each connected profinite groupoid.
The full one-object subcategories have automorphism groups $P,P'$;
they are equivalent to the original groupoids because every other
object is isomorphic to the chosen one. A functor is then described
by a continuous group homomorphism $\varphi:P\to P'$. A natural
transformation to the functor of $\psi$ is precisely an element
$q\in P'$ with $\psi=\operatorname{int}(q)\varphi$, and the
natural automorphisms of the first functor are
$C_{P'}(\varphi(P))$. This proves the rigidity, up to a unique
natural isomorphism, asserted in the Scholie on scans15--16.
The corresponding finite local systems are the finite continuous
$P$-sets, so passage to their classifying topoi gives the same
comparison without changing these natural transformations.
\end{proof}

\begin{result}{B12}{Refuted}{The semiabelian Albanese homology claim for arbitrary reduced sources}
\label{res:B12}
For a general reduced connected finite-type scheme, a universal
semiabelian recipient of pointed morphisms need not have the same
$\ell$-adic $H_1$. Thus the assertion numbered 3 on scan 25 does not
justify Corollary 1(a) in all of its stated source generality.
B3 supplies a proof in that generality. This is sweep row ZAB5
(149, section 1, manuscript p.~12, scan 25).
\end{result}
\begin{proof}
Over $\mathbb C$ take the nodal rational cubic
\[
             X:\quad Y^2Z=X^2(X+Z).
\]
Its normalization is $\mathbb P^1$, with map
\[
 [T:S]\longmapsto
 [S(T^2-S^2):T(T^2-S^2):S^3].
\]
The two points $[1:1],[-1:1]$ map to the node; elsewhere this is
an isomorphism. The displayed equation and the node are checked
in \texttt{code/AB\_checks.py}. A morphism $X\to G$ to any
semiabelian variety becomes constant on its normalization:
the abelian quotient admits no nonconstant map from $\mathbb P^1$,
and the remaining torus-valued map on the proper curve is constant.
Surjectivity of normalization implies the original map is constant.
Thus its universal pointed semiabelian recipient is the zero group.

Topologically, identifying two distinct points of the two-sphere
produces a space homotopy equivalent to $S^2\vee S^1$. Its finite
cyclic covers can equivalently be constructed by taking a cyclic
chain of $n$ copies of $\mathbb P^1$ and identifying successive
preimages of the node. Each node has precisely the same two-branch
local model, so these are finite \'etale covers of $X$.
Complex finite-cover comparison therefore gives
$\pi_1(X)=\widehat{\mathbb Z}$ and
$H_1(X,\mathbb Z_\ell)=\mathbb Z_\ell$.
The unique map to the zero semiabelian recipient induces the zero
map on this nonzero module. This is the exact failed comparison.
\end{proof}

\begin{result}{B13}{Refuted}{The toric-kernel disjointness assertion on scan 38}
\label{res:B13}
Let $Y=C\setminus D$ be a smooth hyperbolic curve over an algebraically
closed field of characteristic zero, and let $J_D\to J_C$ be its
compact-Jacobian quotient. For $g(C)\geq1$, every nonzero
$u\in\ker(J_D\to J_C)$ satisfies
$(\alpha(Y)+u)\cap\alpha(Y)=\varnothing$.
This assertion is false for genus zero. There, for any nonidentity
translation, the intersection has at most one geometric point, and
its existence is determined by the equations below. This corrects
the bracketed claim and the margin ``$g\geq2$?'' on scan 38,
sweep row ZAB8 (149, section 1, manuscript p.~18).
\end{result}
\begin{proof}
For positive genus, projecting $\alpha(s)=\alpha(t)+u$ to $J_C^1$
gives equality of the compact Abel images of $s,t$. This forces
$s=t$: otherwise a principal divisor $s-t$ would give a degree-one
function on $C$, impossible for positive genus. The original equality
then gives $u=0$. This proves the first assertion, with the explicitly
necessary nonzero hypothesis.

For genus zero put three of the missing points at $0,1,\infty$, and
write the other finite missing points as $d_i$. Up to constant
coordinates determined by the choice of Abel origin, the generalized
Abel immersion is
\[
 \alpha(t)=(t,t-1,(t-d_i)_i)\in\mathbb G_m^{r-1}.
\]
Let $u=(a,b,(c_i)_i)$. A point of the intersection is specified exactly
by $t,s\in Y$ satisfying
\[
 s=at,\qquad s-1=b(t-1),\qquad s-d_i=c_i(t-d_i).
\]
If $a=b$, the first two equations give $a=b=1$; then $s=t$ and the
remaining equations give every $c_i=1$, because $t$ is not a cusp.
This is the identity translation. Otherwise
\[
 t=\frac{1-b}{a-b},\qquad s=\frac{a(1-b)}{a-b}.
\]
Thus the intersection has exactly one point precisely when these
values avoid all missing points and satisfy the remaining displayed
equations; it is empty otherwise. For three punctures take
$(a,b)=(2,3)$: then $t=2,s=4$, and
$(2,1)(2,3)=(4,3)$. This is the promised counterexample in the
entire toric kernel $J_D\to J_C=0$. The integer calculation is
included in \texttt{code/AB\_checks.py}. A single intersection
point supplies no translated open subset, so the criterion is
consistent with the genus-zero part of B4.
\end{proof}
